% !PRINTED 566
\weberpage{566}

\begin{center}
\S.\,156.

\gesperrt{Application to algebraic fields.}
\end{center}

From the minimum for a positive quadratic form one can
derive a minimum of linear forms, by conceiving a positive
quadratic form as a sum of squares of linear forms. On this
rest the following applications to the theory of algebraic
fields, which we preface with a lemma:

1. \gesperrt{If $a_1, a_2, \ldots, a_n$ are real positive numerical values, then}
\weq{1}{a_1 + a_2 + \cdots + a_n \geqq n \sqrt[n]{a_1 a_2 \ldots a_n}.}

This theorem is evidently correct for $n=2$, for
$a_1 + a_2 - 2\sqrt{a_1 a_2} = (\sqrt{a_1} - \sqrt{a_2})^2$, and thus never negative. We
therefore apply complete induction, by assuming formula
(1) as already proved for $n-1$ terms. If we then assume
that $a_n$ is exceeded in magnitude by none of the remaining
$a$, then
\[ a_n \geqq \sqrt[n-1]{a_1 a_2 \ldots a_{n-1}} = b, \]
and
\[ a_1 + a_2 + \cdots + a_{n-1} \geqq (n-1)\,b. \]
Hence
\[ a_1 + a_2 + \cdots + a_n \geqq a_n + (n-1)\,b, \]
and consequently
\[ \frac{a_1 + a_2 + \cdots + a_n}{n} \geqq b + \frac{a_n - b}{n}, \]

\begin{footnotesize}
\noindent\rule{4cm}{0.4pt}

in which $\Gamma(x)$ denotes the gamma function, for which the equations
\[ \Gamma(x) = (x-1)\,\Gamma(x-1), \quad \Gamma\left(\frac{1}{2}\right) = \sqrt{\pi} \]
hold, so that $K_n$ is expressed by rational numbers and by
$\sqrt{\pi}$.

(Cf. e.\,g. Vorlesungen \"uber die Theorie der bestimmten Integrale
nach \gesperrt{Lejeune-Dirichlet} von G.\,F. \gesperrt{Meyer} herausgegeben, \S.\,175.
Leipzig 1871).

From this there results
\[ \Theta_n = \frac{4}{n\pi}\,\Gamma\left(1 + \frac{n}{2}\right)^{\frac{2}{n}}, \]
which, by the approximation formulae for the $\Gamma$-function, for large
values of $n$ agrees nearly with
\[ \frac{2}{\pi e} = 0{,}234\ldots \]
(\gesperrt{Minkowski}, Crelle, Bd. 107, S. 292).
\end{footnotesize}
%CONT>>
